\documentclass[11pt]{article}

\usepackage[T1]{fontenc}
\usepackage[utf8]{inputenc}
\usepackage{lmodern}
\usepackage[margin=1in]{geometry}
\usepackage{amsmath,amssymb,amsthm,mathtools,bm}
\usepackage{longtable,booktabs,array,tabularx}
\usepackage[shortlabels]{enumitem}
\usepackage{xcolor}
\usepackage{hyperref}
\usepackage{xurl}
\usepackage{microtype}
\usepackage{seqsplit}

\hypersetup{colorlinks=true,linkcolor=blue,citecolor=blue,urlcolor=blue}
\setlength{\parindent}{0pt}
\setlength{\parskip}{0.55em}
\setlength{\emergencystretch}{3em}
\renewcommand{\arraystretch}{1.12}
\newcolumntype{L}[1]{>{\raggedright\arraybackslash}p{#1}}
\newcolumntype{R}[1]{>{\raggedleft\arraybackslash}p{#1}}

\newtheorem{theorem}{Theorem}
\newtheorem{definition}{Definition}
\newtheorem{proposition}{Proposition}
\theoremstyle{remark}
\newtheorem{remark}{Remark}

\newcommand{\OPH}{\mathrm{OPH}}
\newcommand{\U}{\mathrm{U}}
\newcommand{\SU}{\mathrm{SU}}
\newcommand{\A}{\mathcal A}
\newcommand{\R}{\mathbb R}
\newcommand{\Z}{\mathbb Z}
\newcommand{\dd}{\,\mathrm d}
\newcommand{\e}{\mathrm e}
\newcommand{\alphainv}{\alpha^{-1}}
\newcommand{\ATh}{A_{\mathrm{Th}}}
\newcommand{\AZ}{A_Z}
\newcommand{\Pstar}{P_\star}
\newcommand{\Ppub}{P_{\mathrm C}}
\newcommand{\Rq}{R_Q}
\newcommand{\Dcalc}{\Delta_{\mathrm{calc}}}
\newcommand{\Dcalib}{\Delta_{\mathrm{H,cal}}}
\newcommand{\Dlep}{\Delta_{\mathrm{lep}}}
\newcommand{\Dhad}{\Delta_{\mathrm{had}}}
\newcommand{\DEW}{\Delta_{\mathrm{EW}}}
\newcommand{\Dreq}{\Delta_{\mathrm{req}}}
\newcommand{\Dhadreq}{\Delta_{\mathrm{H,req}}}
\newcommand{\Dmiss}{\Delta_{\mathrm{miss}}}
\newcommand{\Nc}{N_c}
\newcommand{\Ng}{N_g}
\newcommand{\decnum}[1]{\vcenter{\hbox{\parbox{0.68\linewidth}{\small\ttfamily\seqsplit{#1}}}}}

\definecolor{ophblue}{HTML}{1B4E7A}
\definecolor{ophgreen}{HTML}{1F7A5A}
\definecolor{ophgold}{HTML}{A66A00}
\definecolor{ophred}{HTML}{8A2D2D}

\title{\textbf{The Fine-Structure Constant as an OPH Pixel Fixed Point}}
\author{Bernhard M\"uller\\Pragma Inc.
\and
Alexander Osika\\SNRGY Studios}
\date{}
\newcommand{\OPHPaperReleaseID}{r1520}
\newcommand{\OPHPaperReleaseDate}{July 9, 2026}
\newcommand{\OPHPaperReleaseBanner}{%
\begin{center}
\small\textbf{Paper release:} \texttt{\OPHPaperReleaseID}\quad\textbf{Released:} \OPHPaperReleaseDate
\end{center}
}


\begin{document}
\maketitle
\begin{center}
\small\textbf{Paper release:} \OPHPaperReleaseID\quad
\textbf{Released:} \OPHPaperReleaseDate
\end{center}

\begin{abstract}
Observer Patch Holography (OPH) assigns one screen cell two readings: an
outside dimensionless area \(P\) and the electromagnetic observation scale
seen from within. Principle SL-2 of the strange-loop principle set
declares the detuning law
\[
P=\varphi+\frac{\sqrt\pi}{\ATh(P)},
\]
and principle SL-3 identifies the measured fine-structure constant with the
substrate pixel readout, so the working value of \(P\) is located from CODATA
\(\alpha\) pending exact closure; the theory layer takes no measured number.
The source-only
calculation contracts to \(\alpha^{-1}=136.994835177413\ldots\), the
interval-certified unique fixed point of the declared source map (enclosure
width \(7.2\times10^{-24}\)); it sits outside the SL-3 basin by
\(3.0\times10^{-4}\) in relative units, and this loop residual is closure row
CL-1 of the ledger. The self-consistent gauge-width map has the certified
fixed point \(\alpha^{-1}=137.035660136946577\ldots\), \(2.5\times10^{-6}\)
below the measured \(137.035999177(21)\), about \(1.6\times10^{4}\)
measurement sigma (row CL-2); the previously displayed
\(137.0359595008\ldots\) is a mixed-provenance display packet (source root
plus \(\alpha_U\) at the SL-3 pixel), not a fixed point of any single declared
map, and stays on record as a display packet only.
Neither number is a derivation of \(\alpha\): under SL-3 the measured value is
the input that fixes the pixel, and both residuals are attributed to the same
missing term. The paper traces that term to low-energy quark and gluon
transport. Confinement makes a same-scheme, non-perturbative hadronic spectral
calculation necessary; that source calculation is work in progress.
\end{abstract}

\noindent\textbf{Keywords:} fine-structure constant, Observer-Patch Holography, pixel fixed
point, electromagnetic coupling, Standard Model constants, hadronic vacuum polarization

\section*{What This Paper Contributes}

The fine-structure constant is usually treated as an empirical constant. This paper
turns it into a branch readout of one OPH screen cell. The known physics is the
low-energy Maxwell coupling, charged-lepton running, and low-energy hadronic
spectral transport. OPH adds the pixel self-reference: the outside area coordinate
\(P\) must agree with the inside electromagnetic observation scale read by
observers on the same branch.

The pixel map carries zero continuous dials, and each declared map has a
machine-certified unique fixed point: the source map contracts to
\(136.994835177413\ldots\) (\(3.0\times10^{-4}\) relative from the measured
endpoint) and the self-consistent gauge-width map to
\(137.035660136947\) (\(2.5\times10^{-6}\) relative; the declared map is
certified, the physical transport step is open). The remaining
gap has an explicit address: the hadronic transport needed in the same
electromagnetic convention as the source map. This split is the contribution: a fixed-point
equation for the pixel branch with no adjustable form, a visible address for the missing low-energy
QCD step, and a clear boundary between source theorem, empirical endpoint
closure, and measured comparison.

\section{Introduction}

The goal is to make every mathematical step visible and keep the physics boundary clear.

Two declared principles carry the construction; their canonical statements live in
the project's strange-loop principle register and are cited here rather than restated. Principle SL-2
(detuning law) declares that the pixel ratio sits off the golden-ratio balance
point by the Gaussian-normalized width of one electromagnetic observation,
\(P=\varphi+\sqrt\pi\,\alpha\). Principle SL-3 (constant identification) declares
that under self-simulation the fine-structure constant inside the world and
the substrate pixel readout are one quantity; the working value of \(P\) is
located from measured \(\alpha\) pending exact closure, and the theory layer
takes no measured number. Neither principle is derived in this paper; each is a starting point of
deduction, and evidence bears on the pair through the closure-ledger residuals
CL-1 and CL-2 quoted throughout.

In the selection chain of the project consistency stack,
the detuning law is row C6: after C1--C5
have removed the structural freedoms, C6 is the step that leaves exactly one
local freedom, the pixel, and pins it to one equation. That exactness is what
makes the loop residuals CL-1 and CL-2 decisive: a detuning law with an
adjustable form would turn each residual into a tunable band, while the
declared law makes each residual a measurement of the principle set.

There is a long history of attempts to understand \(\alpha\) from deeper structure. Jentschura and Nandori give a useful survey of first-principles attempts, including beta-function, symmetry, cutoff, and string-inspired ideas \cite{jentschura-nandori}. Golden-ratio constructions appear in historical geometric proposals \cite{heyrovska-narayan} and in semi-empirical work on electroweak and flavor mixing \cite{ciborowski}. Octonionic and exceptional-algebra approaches also produce values near \(1/137\) \cite{singh}. ArXiv papers use symbolic regression and quantum-information criteria to look for structure in Standard Model constants and electroweak parameters \cite{chekanov-kjellerstrand,liu-low-yin}. The OPH calculation belongs to this broad search for structure, with one added requirement: the number must arise from a fixed-point map whose input and output are the same local screen cell.

\begin{quote}\small
\textbf{Source/root audit.} The executable source map contracts to the undressed source/root inverse coupling \(\alphainv_{\mathrm{root}}=136.994835177413\ldots\), the interval-certified unique fixed point of the declared map (interval certificate shipped with the code release \cite{oph-code}; enclosure width \(7.2\times10^{-24}\), \(L\le0.0724\)). The earlier printed tail \(136.9948351646\ldots\) came from an unconverged run and is superseded beyond digit~9 (ledger row CL-6, closed). The self-consistent map that adds the finite-screen unified gauge width \(\alpha_U(P)\) inside the loop has the certified fixed point \(137.035660136946577\ldots\); this is the current closure status of the loop (ledger row CL-2, \(2.5\times10^{-6}\) relative, about \(1.6\times10^{4}\) measurement sigma). The historically displayed combination \(A_{\alpha_U}^{\mathrm{fp}}=137.0359595008\ldots\) mixes the source root with \(\alpha_U(\Ppub)\) evaluated at the SL-3 pixel (the CODATA-derived comparison pixel); it is a mixed-provenance display packet, not a fixed point of any single declared map, and stays on record as a display packet only. Neither certified fixed point is a source-only prediction or a derivation of \(\alpha\), because measured \(\alpha\) is the input that fixed the pixel under SL-3.

\textbf{CODATA/NIST measured endpoint.} CODATA/NIST reports \(\alphainv(0)=137.035999177(21)\). The difference from \(A_{\alpha_U}^{\mathrm{fp}}\) is CODATA-minus-diagnostic bookkeeping; it is not the output of a same-scheme hadronic spectral calculation.
\end{quote}

\paragraph*{Intuitive picture: three readings of one pixel.}
The three fine-structure numbers are not three unrelated corrections. They are
three increasingly observer-facing readings of the same local pixel branch:
\[
\text{source pixel}
\longrightarrow
\text{gauge-width dressed pixel}
\longrightarrow
\text{measured low-energy value},
\]
or numerically,
\[
136.994835177413\ldots
\longrightarrow
137.035660136946577\ldots
\longrightarrow
137.035999177(21).
\]
The first two values are the interval-certified fixed points of the declared
source map and of the declared self-consistent gauge-width map
(interval contraction certificate shipped with the code release \cite{oph-code}).
The root value \(\alpha_{\mathrm{root}}^{-1}\) is the naked source/root inverse
coupling: the electromagnetic width emitted when one screen cell closes as a
geometric self-consistency object before the low-energy charged vacuum is
attached. The outside reading is the pixel detuning
\((P-\varphi)/\sqrt\pi\); the inside reading is the electromagnetic observation
strength on the same branch. The fixed point says those two readings must name
one cell.

The contribution \(\alpha_U(\Ppub)\) is the finite-screen unified gauge-width evaluated at the
CODATA-derived comparison pixel. It is the main high-scale gauge-sector width carried by that
pixel: the weak and color representation content must fit the pixel budget
\(P/4\), and the D10 heat-kernel closure supplies the executable certificate.
Adding this comparison-pixel width to the source root gives the mixed-provenance no-hadron
diagnostic \(A_{\alpha_U}^{\mathrm{fp}}\); that additive combination is a display packet,
not a fixed point of any single declared map. The self-consistent reading, with
\(\alpha_U(P)\) evaluated inside the loop, is the certified fixed point
\(137.035660136946577\ldots\) quoted above. Both are close to the laboratory endpoint because
the gauge width supplies almost all of the inverse-alpha gap from the naked source value.

The factor \(C_{24,Q}\) is a back-solved calibration: the small endpoint
dressing required after the gauge-width contribution is added, solved from the
measured endpoint rather than computed from source data. It is near one
because \(\alpha_U(\Ppub)\)
supplies the large part of the correction. Its excess
\(\alpha_U(\Ppub)(C_{24,Q}-1)\) is a comparison residual with the numerical size required of the
remaining same-scheme low-energy hadronic endpoint transport. It does not compute that transport. A zero-momentum laboratory photon does
not probe an empty charged vacuum; it sees the Ward-projected \(\mathrm U(1)_Q\)
current after charged-lepton vacuum polarization, confined-quark/hadron
spectral transport, and finite endpoint matching. In short, the observer
measures the fully dressed low-energy coupling, while
\(A_{\alpha_U}^{\mathrm{fp}}\) is a mixed-provenance no-hadron diagnostic before that dressing.

The OPH background used here is spread across three papers. The observer-overlap starting point is the synthesis paper \cite{oph-synthesis}. The source map and fixed-point branch are recorded in the compact overlap-consistency paper \cite{oph-smgr}. The charged-spectrum and particle-structure continuation is recorded in the particle paper \cite{oph-particle}. The references give direct GitHub links to those source files.

Why this split matters: the fixed-point equation and the hadronic transport calculation answer different questions. The fixed-point equation says how a local screen cell must close on itself. The hadronic calculation says how the electromagnetic current is transported through low-energy QCD. These are different layers of the calculation. A non-arbitrary endpoint calculation must either evaluate the same-scheme hadronic spectral functional directly, or derive the scheme bridge from OPH source data.

\begin{proposition}[Maxwell normalization of the endpoint lane]
On the ordinary electromagnetic branch used by the fixed-point equation, the Ward-projected
\(\U(1)_Q\) channel carries
\[
F_Q=dA_Q,\qquad dF_Q=0,\qquad d{*}F_Q=g_Q^2(q^2;P){*}J_Q,
\]
with
\[
g_Q^2(q^2;P)=4\pi\alpha_{\mathrm{em}}(q^2;P),
\qquad
\ATh(P)=\alpha_{\mathrm{em}}^{-1}(0;P).
\]
Thus the inside reading in the pixel equation is the Thomson-limit Maxwell coupling of the same
electromagnetic current.
\end{proposition}

\begin{proof}
The compact OPH reconstruction identifies the unbroken electromagnetic branch as \(\U(1)_Q\).
On that abelian factor, the compact-gauge curvature reduces to \(F_Q=dA_Q\), and the quadratic
field action gives \(d{*}F_Q=g_Q^2{*}J_Q\). The particle-paper transport theorem then identifies
the same \(g_Q(q^2;P)\) with the Ward-projected running electromagnetic coupling. Taking
\(q^2\to0\) gives the endpoint quantity \(\ATh(P)\) used in the pixel fixed-point equation.
\end{proof}

The hadronic step uses the standard dispersion logic behind hadronic vacuum-polarization work. The relevant comparison data are electromagnetic spectral functions measured through \(e^+e^-\to\mathrm{hadrons}\), as used in data-driven running-\(\alpha\) and \(g-2\) analyses \cite{dhmz,knt19,pdg-xsect}. In this paper, the decimal correction shown above is reported as the difference between the CODATA/NIST comparison value and the OPH pure source calculation. A direct OPH production calculation of the hadronic spectral function remains open.

\clearpage
\section{Symbols}

\begin{longtable}{@{}L{0.20\textwidth}L{0.32\textwidth}L{0.38\textwidth}@{}}
\toprule
Symbol & Definition & Informal meaning \\
\midrule
\endhead
\(\A(P)\) & Patch algebra assigned to a screen patch \(P\) & The local observables available to one finite observer patch. \\
\(\omega_P\) & State on \(\A(P)\) & The local physical description carried by that patch. \\
\(\A(P\cap Q)\) & Shared overlap algebra & The observables two neighboring patches can compare. \\
\(P\) & \(a_{\mathrm{cell}}/\ell_\star^2\) & The dimensionless area of one screen cell in scale-readout units. \\
\(a_{\mathrm{cell}}\) & Physical screen-cell area & The area of the local pixel on the holographic screen. \\
\(\ell_\star\) & OPH scale-certificate length & The length emitted by the selected no-\(G\) scale certificate. \\
\(\ell_P\) & Planck length & The SI display of \(\ell_\star\) after \(G_{\mathrm{SI}}=c^3\ell_\star^2/\hbar\). \\
\(\varphi\) & \((1+\sqrt5)/2\) & The self-similar entropy-balance point. \\
\(\sqrt\pi\) & Boundary Gaussian normalization width & The conversion factor from pixel displacement to observation width. \\
\(\alpha_{\mathrm{ext}}(P)\) & \((P-\varphi)/\sqrt\pi\) & The outside, geometric reading of the electromagnetic coupling. \\
\(\ATh(P)\) & \(\alpha_{\mathrm{em}}^{-1}(0;P)\) & The inverse electromagnetic coupling at the Thomson limit emitted by the trial pixel. \\
\(\alpha_{\mathrm{in}}(P)\) & \(1/\ATh(P)\) & The inside, electromagnetic reading of the same cell. \\
\(E_P\) & Planck energy & The energy unit associated with the selected scale certificate. The numerical equations use Planck units unless GeV labels are attached. \\
\(M_U(P)\) & \(E_P\e^{-2\pi}P^{1/6}\) & The OPH unification-scale readout emitted by a trial pixel. \\
\(E_{\mathrm{cell}}(P)\) & \(E_P/\sqrt P\) & The local cell energy readout. \\
\(\alpha_U(P)\) & Unified coupling solved from the heat-kernel closure equation & The coupling at \(M_U(P)\) that makes the gauge representation entropy match the pixel. \\
\(\alpha_i(\mu;P)\) & Gauge coupling \(i=1,2,3\) at scale \(\mu\) & The hypercharge, weak, and color couplings emitted by the same trial \(P\). \\
\(b_i\) & \((33/5,1,-3)\) & The one-loop beta coefficients for the high-scale running convention used here. \\
\(m_Z(P)\) & Self-consistent electroweak anchor scale & The \(Z\)-scale generated by \(v(P)\), \(\alpha_1(P)\), and \(\alpha_2(P)\). \\
\(v(P)\) & \(E_{\mathrm{cell}}(P)\exp[-2\pi/(\beta_{\mathrm{EW}}\alpha_U(P))]\) & The electroweak transmutation scale. \\
\(\beta_{\mathrm{EW}}\) & \(\Nc+1=4\) for \(\Nc=3\) & The coefficient that controls the exponential drop to the electroweak scale. \\
\(\AZ(P)\) & \(\alpha_{\mathrm{em}}^{-1}(m_Z^2;P)\) & The electroweak-scale electromagnetic anchor. \\
\(\Dlep(P)\) & Lepton transport contribution to inverse alpha & The exact one-loop kernel evaluated on the frozen phenomenological \(e,\mu,\tau\) spectrum. \\
\(\Dhad(P)\) & Hadronic transport contribution & The low-energy QCD spectral contribution needed for the pure source calculation. \\
\(\DEW(P)\) & Electroweak/scheme endpoint remainder & Same-scheme finite remainder needed to connect the source anchor to the endpoint. \\
\(\Dcalc(P)\) & Calculated charged-fermion transport contribution & The exact kernel on the supplied lepton continuation plus the screened quark continuation; not a pure source prediction. \\
\(\Dcalib^{\mathrm{fp}}\) & Root-to-CODATA bookkeeping difference & The inverse-alpha amount between the naked root value and the CODATA/NIST central endpoint; it includes the comparison-pixel \(\alpha_U(\Ppub)\) contribution and is not the post-\(\alpha_U\) hadronic gap. \\
\(\Dhadreq^{\mathrm{fp}}\) & Required hadronic endpoint correction after \(A_{\alpha_U}^{\mathrm{fp}}\) & The inverse-alpha amount by which the displayed OPH output misses the CODATA/NIST central endpoint. \\
\(\Rq(P)\) & \(\ATh(P)-[\AZ(P)+\Dcalc(P)]\) & The remaining same-scheme hadronic endpoint contribution. \\
\(A_{\alpha_U}^{\mathrm{fp}}\) & \(\alpha_{\mathrm{root}}^{-1}+\alpha_U(\Ppub)\) & Mixed-provenance no-hadron comparison diagnostic; the root is source-side and \(\Ppub\) is CODATA-derived. \\
\(J_{24,Q}(P),\omega_Q(P),\Xi_Q(P)\) & Source-side hadronic spectral objects required for a final endpoint theorem & The Jacobi/Stieltjes spectral payload, normalization, and finite same-scheme remainder that must be emitted without looking at the Thomson target. \\
\(\Ppub\) & CODATA/NIST comparison pixel & The pixel obtained from the CODATA/NIST Thomson endpoint under the OPH outer equation. \\
\bottomrule
\end{longtable}

\section{Shared Collar Convention for \(\chi_\nu\)}

The \(\chi_\nu\) protected-reserve theorem uses the public endpoint branch
convention
\[
P_\chi=P_{\mathrm C}=1.630968209403959\ldots .
\]
Its scalar reserve density is
\[
\epsilon_{\rm res}=\frac{P_\chi}{24}
\]
only after the same-collar shared-edge budget and one-class \(\mathbb Z_6\)
trace receipts pass in the screen microphysics and susceptibility papers.  This
does not make \(P/4\) a primitive Hilbert-space dimension.  In this paper,
\(P/4\) is the local screen-cell entropy budget used by the electromagnetic
fixed-point lane; \(P_\chi/24\) is the shared scalar-reserve density used by the
\(\chi_\nu\) collar branch.

\section{Step 1: The OPH Starting Rule}

OPH starts from a finite screen covered by observer patches. If \(P_1\) and \(P_2\) overlap, then their induced states must agree on the shared algebra:
\begin{equation}
\omega_{P_1}|_{\A(P_1\cap P_2)}
=
\omega_{P_2}|_{\A(P_1\cap P_2)}.
\label{eq:overlap}
\end{equation}

Informally: no observer sees the whole world. Physics is what survives when neighboring local descriptions can be checked against each other and made consistent on the parts they share.

The quantitative fine-structure branch applies this rule to a single local screen cell. The cell has one outside description and one inside description. The outside description is geometric. The inside description is electromagnetic. Closure means both descriptions identify the same cell.

Why this step is needed: OPH does not start by assigning constants to nature. It starts by asking what finite observers can compare. A dimensionless constant can be derived only if it is the fixed value of such a comparison.

\section{Step 2: The Pixel Variable}

Define the local pixel ratio
\begin{equation}
P:=\frac{a_{\mathrm{cell}}}{\ell_\star^2}.
\label{eq:Pdef}
\end{equation}
Here \(a_{\mathrm{cell}}\) is the screen-cell area and \(\ell_\star^2\) is the area scale emitted
by the separate OPH gravity readback fixed point.

Informally: \(P\) says how large the screen cell is when measured in scale-readout units. It is dimensionless, so it can be compared directly with pure numbers such as \(\varphi\) and \(\sqrt\pi\). After the gravity row emits \(G\), the same scale is displayed as the Planck area in SI units. The pixel equation fixes this dimensionless ratio; it does not determine the SI scale by itself.

Why this step is needed: \(\alpha\) has no units. A unitless electromagnetic number has to be matched to a unitless geometric number. Dividing the cell area by the emitted scale area gives that number without using the pixel equation to derive \(G\).

\section{Step 3: The Golden-Ratio Balance}

The self-similar balance point is
\begin{equation}
\varphi=1+\frac{1}{\varphi},
\qquad
\varphi^2-\varphi-1=0,
\qquad
\varphi=\frac{1+\sqrt5}{2}.
\label{eq:phi}
\end{equation}

Informally: \(\varphi\) is the fixed point of the simplest self-similar split, where the whole-to-large ratio equals the large-to-small ratio. OPH takes it as the zero-detuning balance point for the local pixel by principle SL-2; the choice is declared, not deduced from prior structure.

The pixel does not sit exactly at \(\varphi\). It sits at
\begin{equation}
\Delta_P:=P-\varphi.
\label{eq:deltaP}
\end{equation}

Informally: \(\Delta_P\) is the small amount by which the realized cell is displaced from exact self-similar equilibrium. That small displacement is what becomes the electromagnetic observation strength.

Why this step is needed: a fixed point needs a reference position. Principle SL-2 declares the golden ratio as the self-similar reference, while \(P-\varphi\) measures how far the actual cell sits away from it.

\section{Step 4: Boundary Gaussian Normalization}

The boundary normalization converts \(\Delta_P\) into the outside coupling readout:
\begin{equation}
\alpha_{\mathrm{ext}}(P):=\frac{P-\varphi}{\sqrt\pi}.
\label{eq:alphaext}
\end{equation}

Informally: the screen displacement is not itself the electromagnetic coupling. The boundary Gaussian width supplies the normalization. Dividing by \(\sqrt\pi\) turns the geometric displacement into a dimensionless observation strength. The Gaussian choice is the second half of the SL-2 declaration: declared principle, not deduction.

Why this step is needed: the outside screen variable is an area displacement. The inside electromagnetic variable is a coupling. The boundary normalization is the conversion factor between those two readings.

\section{Step 5: The Inside Electromagnetic Readout}

Let
\begin{equation}
\ATh(P):=\alpha_{\mathrm{em}}^{-1}(0;P)
\label{eq:ATdef}
\end{equation}
be the inverse electromagnetic coupling at zero momentum, emitted from the same trial pixel \(P\). The inside coupling is
\begin{equation}
\alpha_{\mathrm{in}}(P):=\frac{1}{\ATh(P)}.
\label{eq:alphain}
\end{equation}

Informally: an observer inside the encoded world does not see \(P\) as a screen area. The observer sees the strength of electromagnetism. Since particle physicists quote the low-energy coupling as an inverse number near \(137\), the actual coupling is \(1/\ATh(P)\).

Why this step is needed: the measured fine-structure constant is the Thomson-limit coupling. The calculation must therefore end at zero momentum, after the electroweak anchor has been transported through all charged degrees of freedom.

\section{Step 6: The Fixed-Point Equation}

Closure requires the outside and inside couplings to agree:
\begin{equation}
\alpha_{\mathrm{ext}}(P)=\alpha_{\mathrm{in}}(P).
\label{eq:couplingmatch}
\end{equation}
Using Eqs.~\eqref{eq:alphaext} and \eqref{eq:alphain},
\begin{equation}
\frac{P-\varphi}{\sqrt\pi}
=
\frac{1}{\ATh(P)}.
\end{equation}
Equivalently,
\begin{equation}
\boxed{
H(P):=P-\varphi-\frac{\sqrt\pi}{\ATh(P)}=0
}
\label{eq:H}
\end{equation}
or
\begin{equation}
\boxed{
P=G(P):=\varphi+\frac{\sqrt\pi}{\ATh(P)}.
}
\label{eq:G}
\end{equation}

Informally: feed a trial pixel into the source chain. It emits a Thomson endpoint. That endpoint tells the pixel what size it should have. The physical pixel is the value that comes back unchanged.

Why this step is needed: a trial value is not enough. The same cell must agree with itself when read from the outside and from the inside. That is why the answer is a fixed point, not a one-way evaluation.

\section{Step 7: The Source Map}

For a trial \(P\), the source map first emits
\begin{align}
M_U(P)&=E_P\,\e^{-2\pi}\,P^{1/6},
\label{eq:MU}\\
E_{\mathrm{cell}}(P)&=\frac{E_P}{\sqrt P}.
\label{eq:Ecell}
\end{align}

Informally: \(M_U(P)\) is the high-scale source readout, and \(E_{\mathrm{cell}}(P)\) is the local cell energy. Both come from the same screen cell. The unification scale is therefore read from the pixel.

Why this step is needed: the inside electromagnetic coupling has to be generated from the same \(P\) used in the outside equation. These two formulas begin that inside calculation without adding a separate high-scale fit parameter.

\section{Step 8: Electroweak Transmutation}

The source branch uses
\begin{equation}
\beta_{\mathrm{EW}}=\Nc+1=4
\end{equation}
with \(\Nc=3\), and defines
\begin{equation}
v(P,\alpha_U)
=
E_{\mathrm{cell}}(P)
\exp\!\left[-\frac{2\pi}{\beta_{\mathrm{EW}}\alpha_U}\right].
\label{eq:v}
\end{equation}

Informally: the weak scale is exponentially lower than the cell scale. The unified coupling controls that descent, so changing \(\alpha_U\) changes the electroweak scale sharply.

Why this step is needed: the fine-structure constant is measured at low energy, but the source map begins at the cell scale. The transmutation formula explains how the electroweak scale is produced from the same source data. The capacity side of the same pixel chain is carried by the coupling theorem G2-GAP-1, a conditional theorem proven modulo three declared premises CP-1 to CP-3 that places the record capacity at \(\pi\exp[6\pi/(P\alpha_U)]\), certified by interval arithmetic at \(3.53\times10^{122}\).

\section{Step 9: One-Loop Gauge Running}

For \(i=1,2,3\), define
\begin{equation}
(b_1,b_2,b_3)=\left(\frac{33}{5},\,1,\,-3\right),
\label{eq:b}
\end{equation}
and run the couplings by
\begin{equation}
\alpha_i^{-1}(\mu;P,\alpha_U)
=
\alpha_U^{-1}
+ 
\frac{b_i}{2\pi}\log\!\left(\frac{M_U(P)}{\mu}\right).
\label{eq:running}
\end{equation}

Informally: once \(\alpha_U\) is guessed, all three gauge couplings at lower scales are fixed. The same high-scale coupling determines \(\alpha_1\), \(\alpha_2\), and \(\alpha_3\).

Why this step is needed: the electromagnetic coupling is a mixture of the weak and hypercharge couplings. Running the gauge couplings supplies the ingredients needed to form that mixture at the \(Z\)-scale.

\section{Step 10: The Self-Consistent \(Z\)-Scale}

The hypercharge coupling in Standard Model normalization is
\begin{equation}
\alpha_Y(\mu;P)=\frac{3}{5}\alpha_1(\mu;P).
\label{eq:alphaY}
\end{equation}
The tree-level \(Z\)-mass readout is
\begin{equation}
m_Z(\mu;P,\alpha_U)
=
\frac{v(P,\alpha_U)}{2}
\sqrt{4\pi\alpha_2(\mu;P,\alpha_U)+4\pi\alpha_Y(\mu;P,\alpha_U)}.
\label{eq:mZtree}
\end{equation}
The source point uses the self-consistency condition
\begin{equation}
\mu=m_Z(\mu;P,\alpha_U).
\label{eq:mZfixed}
\end{equation}

Informally: the \(Z\)-scale is determined inside the calculation. It is the scale at which the running couplings and the electroweak transmutation formula reproduce their own \(Z\)-mass readout.

Why this step is needed: using a measured \(m_Z\) here would hide experimental input inside the source map. The self-consistency condition keeps the source calculation closed.

\section{Step 11: Heat-Kernel Gauge Closure}

For \(\SU(2)\), the irreducible representations are labeled by \(j=n/2\), \(n=0,1,\ldots,N_2\), with
\begin{equation}
d_j=2j+1,
\qquad
C_2(j)=j(j+1).
\label{eq:su2rep}
\end{equation}
For \(\SU(3)\), representations are labeled by highest weights \((p,q)\), \(0\le p,q\le N_3\), with
\begin{align}
d_{p,q}&=\frac{(p+1)(q+1)(p+q+2)}{2},
\label{eq:su3dim}\\
C_2(p,q)&=\frac{p^2+q^2+pq+3p+3q}{3}.
\label{eq:su3casimir}
\end{align}

For a compact group \(G\) in this finite representation cutoff, define
\begin{align}
Z_G(t)&=\sum_R d_R\,\e^{-tC_2(R)},\\
\bar\ell_G(t)&=
\frac{1}{Z_G(t)}
\sum_R d_R\,\e^{-tC_2(R)}\log d_R.
\label{eq:ellbar}
\end{align}

The heat-kernel source parameters are
\begin{equation}
t_2=4\pi^2\alpha_2(m_Z;P,\alpha_U),
\qquad
t_3=4\pi^2\alpha_3(m_Z;P,\alpha_U).
\label{eq:t23}
\end{equation}
The pixel-closure equation is
\begin{equation}
\boxed{
\bar\ell_{\SU(2)}(t_2)+\bar\ell_{\SU(3)}(t_3)=\frac{P}{4}.
}
\label{eq:pixelclosure}
\end{equation}
For a trial \(P\), Eq.~\eqref{eq:pixelclosure} is solved for \(\alpha_U(P)\).

Informally: the representation entropy carried by the weak and color sectors must match the pixel capacity assigned to the cell. This step is what locks the unified coupling to the same \(P\) that appears in the outer equation.

Why this step is needed: \(\alpha_U\) should not be chosen by hand. The heat-kernel closure turns the finite representation content of the gauge sectors into an equation for \(\alpha_U(P)\).

\section{Step 12: The Electroweak Anchor}

Once \(\alpha_U(P)\) and \(m_Z(P)\) are solved, define
\begin{align}
\alpha_Y(m_Z;P)&=\frac{3}{5}\alpha_1(m_Z;P),\\
\alpha_{\mathrm{em}}(m_Z^2;P)
&=
\left(
\frac{1}{\alpha_2(m_Z;P)}
+\frac{1}{\alpha_Y(m_Z;P)}
\right)^{-1}.
\label{eq:alphaem}
\end{align}
The source-locked electroweak anchor is
\begin{equation}
\boxed{
\AZ(P):=\alpha_{\mathrm{em}}^{-1}(m_Z^2;P).
}
\label{eq:AZ}
\end{equation}
The weak mixing readout is
\begin{equation}
\sin^2\theta_W(m_Z;P)=
\frac{\alpha_{\mathrm{em}}(m_Z^2;P)}{\alpha_2(m_Z;P)}.
\label{eq:sin2}
\end{equation}

Informally: the source map has reached the electromagnetic coupling at the electroweak anchor scale. The low-energy \(1/137\) number appears only after this anchor is transported to zero momentum.

Why this step is needed: \(\AZ(P)\) is the clean place where the source map meets standard electroweak physics. From here the problem becomes a transport problem for the electromagnetic current.

\section{Step 13: Frozen Charged-Spectrum Continuation Used by the Transport}

The one-loop kernel below is exact for supplied fermion masses, but the charged
spectrum supplied to it in this section is a phenomenological continuation,
not an OPH-derived mass theorem.  It uses \(\Nc=3\), \(\Ng=3\), and the
additional assumptions
\begin{equation}
\epsilon=\frac16,
\qquad
\delta=\frac{\beta_{\mathrm{EW}}}{2\Nc\Ng}
=\frac{2}{9}.
\label{eq:epsdelta}
\end{equation}
The physical charged-spectrum derivation is work in progress. It requires
a source-to-mass attachment, a phase-transport law fixing \(\delta=2/9\),
and an attachment of the ordered roots to physical charged-family lines. For a
positive-semidefinite square-root-mass carrier
\[
C=aI+\rho(e^{i\delta}R+e^{-i\delta}R^2),
\qquad
Q=\frac{1+2(\rho/a)^2}{3},
\]
so \(\rho/a=1/\sqrt2\) is exactly equivalent to the empirical Koide value
\(Q=2/3\). At balance this identity is physical throughout the positive
chamber \(|\delta|\leq\pi/12\pmod{2\pi/3}\), and it does not select the phase.

The finite balance normalization has a closed GNS construction. On
\(\mathcal V_K=\mathbf1\oplus\chi\oplus\bar\chi\) with a minimal two-state
orientation record, take the connected algebra
\(B(\mathcal V_K\otimes\mathbb C^2)\simeq M_6(\mathbb C)\) and event
\(E_+=P_0\otimes I_2+P_c\otimes q_+\). Its singlet and oriented charged blocks
both have rank two. Born--L\"uders conditioning of \(I_6/6\) gives
\(p_0=p_c=1/2\), or \(S_c-S_0=\ln2\). In the tracial GNS space,
\(e_0,e_+,e_-\mapsto I,R,R^2\) is unitary, and the canonical square-root
amplitude gives
\[
 a=\sqrt{p_0},\qquad \rho=\sqrt{p_c/2},\qquad
 \frac{\rho}{a}=\frac1{\sqrt2}.
\]
The connected register fixes the normalization left free by the direct-sum
MaxEnt model. Applying Minimal Admissible Realization (MAR) to this
response-local register is an explicit hypothesis. Physical promotion requires
a chiral regular-\(C_3\) carrier and faithful trace-preserving unital
completely positive maps \(\Phi,\Psi\) between the source and physical
response algebras, with \(\Psi\Phi=\mathrm{id}\). If they intertwine the
block records and identify the physical GNS vector with
\(C=(Y_e^\dagger Y_e)^{1/4}\), Kadison--Schwarz makes \(\Phi\) a
\(*\)-monomorphism and an exact \(L^2\) isometry. It therefore carries the
equal source block powers to the physical response. For an accepted finite
chiral checkpoint with three left and three right family modes, positive
kinetic metrics, and a committed neutral Higgs direction,
\(\Phi=\operatorname{Ad}_{J_L\oplus J_E}\). With \(M_F=X_F^2\), this gives
\[
\widehat Y_e=\frac{\sqrt2}{v}J_LM_FJ_E^\dagger,
\qquad
\mathcal M_L=J_LM_FJ_L^\dagger.
\]
The extended branch \(\mathrm{OPH}^{+}_{\rm ch}\) adds graded physical
completion and quotient source-law selection. Given an exhaustive MAR carrier
class with a positive winner gap, a source-closed BV/BRST continuum, and an
interval or contraction certificate for the charged QFT self-map, it selects
one dressed mass readout. A balanced \(C_3\) fixed point gives \(Q=2/3\); a
\(C_3\)-symmetric fixed point with charged attenuation \(\chi_\star\) gives
\[
Q=\frac{1+e^{-2\chi_\star}}{3}.
\]
An off-plane response requires the full response operator. These are
additional branch conditions beyond OPH5. No \(12/24\) receipt supplies that
physical attachment.
Define three Koide roots
\begin{equation}
r_k=1+\sqrt2\cos\!\left(\delta+\frac{2\pi k}{3}\right),
\qquad k=0,1,2,
\label{eq:koide}
\end{equation}
then sort them in increasing order and write the sorted list as
\((r_1,r_2,r_3)\).

The quark exponent vectors are
\begin{equation}
\bm n_u=(2\Nc,\Nc,0)=(6,3,0),
\qquad
\bm n_d=(2\Nc,\Nc+1,\Nc-1)=(6,4,2).
\label{eq:quarkexponents}
\end{equation}
With \(v=v(P)\),
\begin{align}
m_u&=\frac{v}{\sqrt2}\epsilon^6,
&
m_c&=\frac{v}{\sqrt2}\epsilon^3,
&
m_t&=\frac{v}{\sqrt2},\\
m_d&=\frac{v}{\sqrt2}\epsilon^6,
&
m_s&=\frac{v}{\sqrt2}\epsilon^4,
&
m_b&=\frac{v}{\sqrt2}\epsilon^2.
\label{eq:quarkmasses}
\end{align}

For charged leptons the exponent vector is
\begin{equation}
\bm n_e=(7,4,3).
\label{eq:lepexp}
\end{equation}
This vector and the normalization below belong to the same frozen Stage-5
continuation.  In particular, the present corpus does not derive
\(\bm n_e=(7,4,3)\) or the resulting determinant identity
\(\det M_e=v(P)^3/(2\,6^{14})\) from the OPH axioms.  The historical scan fixes
only exponent differences: every common shift \((7+k,4+k,3+k)\) has the same
ratio residual while changing the determinant.  The extra factor \(2^{1/6}\)
below is asserted rather than derived.
Define
\begin{equation}
\log g_c
=
\frac{1}{3}
\sum_{a=1}^3
\log\!\left(
\frac{r_a^2\sqrt2\,6^{n_{e,a}}}{v}
\right),
\qquad
s_0=\e^{-\log g_c},
\qquad
s_e=s_0\,2^{1/6}.
\label{eq:lepscale}
\end{equation}
Then
\begin{equation}
m_e=s_e r_1^2,
\qquad
m_\mu=s_e r_2^2,
\qquad
m_\tau=s_e r_3^2.
\label{eq:lepmasses}
\end{equation}

On the frozen legacy D10 running-tree branch \(P=1.63094\),
\(v=246.76711732749683\,\mathrm{GeV}\), this construction evaluates to
\begin{equation}
\begin{aligned}
m_e^{\mathrm{S5}}&=0.510882243295\,\mathrm{MeV},\\
m_\mu^{\mathrm{S5}}&=105.635282871\,\mathrm{MeV},\\
m_\tau^{\mathrm{S5}}&=1776.579124017\,\mathrm{MeV}.
\end{aligned}
\label{eq:lepmassesstagefive}
\end{equation}
Against the comparison packet the residuals are approximately \(-228\),
\(-219\), and \(-197\) ppm, respectively, with maximum absolute mass error
\(0.02284\%\).  On the public-pixel probe
\(P=1.6309681897\), \(v=246.6174823334856\,\mathrm{GeV}\), the same formula gives
\((0.510572453797,105.571227599,1775.501839453)\,\mathrm{MeV}\) and maximum
error \(0.08347\%\).  These numbers therefore document an accurate, but
branch-sensitive, complete postdiction.  They are not a prospective OPH
prediction: reference masses are absent at runtime, but the frozen formula
family has historical post-hoc ancestry.  Moreover, no charged RG/threshold
theorem maps this historical mass coordinate to the pole masses used for the
ppm comparison. Its executable receipt is committed with the repository
sources.

Informally: this section specifies a frozen continuation spectrum used by the
transport diagnostic.  The lepton vacuum-polarization kernel is an exact
one-loop calculation conditional on those masses; that exactness does not
promote the masses themselves.  The quark part is also a perturbative
continuation. The confined hadronic QCD spectral measure is a separate
low-energy object.

Why this step is needed: vacuum polarization depends on charged particles. The
transport diagnostic cannot be evaluated until a charged spectrum and its
charges are supplied.  A pure OPH source theorem would first have to replace
the continuation spectrum by independently derived charged masses.

\section{Step 14: Exact One-Loop Fermion Transport}

For a fermion of mass \(m_f\), electric charge \(Q_f\), and multiplicity \(N_f\), define
\begin{equation}
K_f(Q^2;m_f,Q_f,N_f)
=
\frac{2N_fQ_f^2}{\pi}
\int_0^1 x(1-x)
\log\!\left(1+\frac{Q^2x(1-x)}{m_f^2}\right)\dd x.
\label{eq:kernelint}
\end{equation}

The integral has the following closed form. Let
\begin{equation}
z=\frac{Q^2}{m_f^2},
\qquad
a=\frac{z}{4}.
\end{equation}
Then
\begin{equation}
\int_0^1 x(1-x)\log(1+zx(1-x))\dd x
=
-\frac{5}{18}
+\frac{1}{6a}
+\frac{(2a-1)\sqrt{1+a}\,\operatorname{asinh}(\sqrt a)}
{6a^{3/2}},
\label{eq:kernelclosed}
\end{equation}
with
\begin{equation}
\operatorname{asinh}(\sqrt a)
=
\log(\sqrt a+\sqrt{1+a}).
\end{equation}

The lepton contribution is
\begin{equation}
\Dlep(P)=
K_e(m_Z(P)^2;m_e,1,1)
+K_\mu(m_Z(P)^2;m_\mu,1,1)
+K_\tau(m_Z(P)^2;m_\tau,1,1).
\label{eq:Dlep}
\end{equation}
The naive five-quark contribution is
\begin{align}
\Delta_q^{\mathrm{naive}}(P)
&=
K_u(m_Z(P)^2;m_u,2/3,3)
+K_d(m_Z(P)^2;m_d,-1/3,3)
\nonumber\\
&\quad
+K_s(m_Z(P)^2;m_s,-1/3,3)
+K_c(m_Z(P)^2;m_c,2/3,3)
+K_b(m_Z(P)^2;m_b,-1/3,3).
\label{eq:Dqnaive}
\end{align}
The perturbative screening factor is
\begin{equation}
S_{\mathrm{calc}}(P)=
1-\frac{\Nc\alpha_3(m_Z;P)}{\pi}.
\label{eq:Scalc}
\end{equation}
Thus
\begin{equation}
\Dcalc(P)
=
\Dlep(P)+S_{\mathrm{calc}}(P)\Delta_q^{\mathrm{naive}}(P).
\label{eq:Dcalc}
\end{equation}

Informally: the supplied charged-lepton continuation is propagated through an
exact one-loop kernel. Quarks are treated by a perturbative expression with a
simple screening factor. The expected gap is the confined hadronic spectral
transport plus the matching terms needed to use the same electromagnetic-current convention all the way to the endpoint.

Why this step is needed: the Thomson endpoint is not the same as the \(Z\)-scale anchor. Charged particles screen the electromagnetic current between those scales, and the kernel is the mathematical form of that screening.

\section{Step 15: The Source Calculation Endpoint}

The calculated endpoint is
\begin{equation}
A_{\mathrm{calc}}(P)
=
\AZ(P)+\Dcalc(P).
\label{eq:Acalc}
\end{equation}
Putting \(A_{\mathrm{calc}}\) into Eq.~\eqref{eq:G} gives the source map
\begin{equation}
G_{\mathrm{calc}}(P)
=
\varphi+\frac{\sqrt\pi}{A_{\mathrm{calc}}(P)}.
\label{eq:Gcalc}
\end{equation}
The numerical solve uses
\begin{equation}
G_{\mathrm{calc}}(P_{\mathrm{root}})=P_{\mathrm{root}}
\label{eq:Proot}
\end{equation}
and emits
\begin{align}
P_{\mathrm{root}}
&=
1.630972095858897376964513903506955628479\ldots,
\label{eq:Prootnum}\\
\alphainv_{\mathrm{root}}
&=
136.9948351774129372952894294644369028576\ldots.
\label{eq:alpharoot}
\end{align}
These are the converged values: the fixed point is interval-certified as
existing and unique on its interval (interval certificate,
shipped with the code release \cite{oph-code};
enclosure widths \(6.8\times10^{-28}\) in \(P\) and \(7.2\times10^{-24}\) in
\(\alpha^{-1}\), Lipschitz bound \(L\le0.0724\)). Earlier releases printed
\(P_{\mathrm{root}}=1.630972095694329\ldots\) and
\(\alphainv_{\mathrm{root}}=136.994835164622\ldots\); those tails came from an
unconverged run and are superseded beyond digit \(\sim\)10 and digit~9
respectively (ledger row CL-6, closed).
The source anchor at the certified root point is
\begin{equation}
\AZ(P_{\mathrm{root}})
=
128.3082680579875973479040576140847588750331\ldots.
\label{eq:AZroot}
\end{equation}
The calculated transport contribution is
\begin{equation}
\Dcalc(P_{\mathrm{root}})
=
8.68656711942533994738537185035214398252\ldots.
\label{eq:Dcalcroot}
\end{equation}

Informally: the fixed-point algebra lands at a stable value near \(137\), with no measured alpha inserted into the solve. The remaining gap has a precise address: low-energy hadronic transport in the same endpoint convention.

Why this step is needed: this is the non-circular check. It shows what the OPH source chain produces before the calibrated hadronic endpoint correction is added.

\section{Step 16: The CODATA/NIST Measured Endpoint and Comparison Pixel}

The CODATA/NIST 2022 inverse fine-structure constant is
\begin{equation}
A_{\mathrm C}=137.035999177,
\qquad
\sigma_A=0.000000021,
\label{eq:codata}
\end{equation}
with concise form \(137.035999177(21)\). The corresponding comparison pixel is
\begin{equation}
\Ppub
=
\varphi+\frac{\sqrt\pi}{A_{\mathrm C}}
=
\decnum{1.6309682094039593248792798477826489413359828516279250606661507533907793398933432}
\label{eq:Ppub}
\end{equation}
The CODATA/NIST central coupling is
\begin{equation}
\alpha(0)
=
\frac{1}{A_{\mathrm C}}
=
\decnum{0.0072973525643314250302457952646916832280660213133653604957798803819933561573639928}
\label{eq:alphapub}
\end{equation}

Informally: once the measured Thomson endpoint is supplied, the outer OPH equation fixes the corresponding comparison pixel immediately. It is the direct readout of one equation.

Why this step is needed: the measured endpoint and the comparison pixel are two forms of the same fixed-point statement. Reporting both makes it clear how the measured inverse coupling maps back to the screen cell.

\paragraph{Two-\(P\) provenance audit.}
There are two distinct pixel numerals in this paper.  The source-only
solve gives \(P_{\mathrm{root}}\) in Eq.~\eqref{eq:Prootnum}; it is the
branch used to audit what the OPH source chain emits before a same-scheme
hadronic endpoint correction is supplied.  The comparison value \(\Ppub\) in
Eq.~\eqref{eq:Ppub} is different by about \(3.9\times10^{-6}\) in \(P\)-space
because it is defined from the measured CODATA/NIST Thomson endpoint.  Thus
the published comparison pixel is a measured-endpoint display, not a
derivation of \(\alpha\).  Downstream quantities that use \(\Ppub\) inherit
that measured endpoint by definition; downstream quantities that use
\(P_{\mathrm{root}}\) are source-branch diagnostics.

\section{Step 17: Endpoint Accounting at the CODATA/NIST Comparison Pixel}

At \(P=\Ppub\), the source point gives
\begin{align}
\AZ(\Ppub)
&=
\decnum{128.30796547328624820996110874175671618724547618036535646005342169635117784168285644078724728}
\label{eq:AZpub}\\
\Dcalc(\Ppub)
&=
\decnum{8.6865678427085284009854425428859697682672217376487364233784577389993784459106783080961179590}
\label{eq:Dcalcpub}
\end{align}
The transport required by the CODATA/NIST measured endpoint is
\begin{equation}
\Dreq(\Ppub)
=
A_{\mathrm C}-\AZ(\Ppub)
=
\decnum{8.72803370371375179003889125824328381275452381963464353994657830364882215831714355921275272}
\label{eq:Dreqpub}
\end{equation}
Therefore the source-side residual at the CODATA/NIST comparison pixel is
\begin{align}
\Rq(\Ppub)
&=
\Dreq(\Ppub)-\Dcalc(\Ppub)
\nonumber\\
&=
\decnum{0.04146586100522338905344871535731404448730208198590711656812056464944371240646525111663476}
\label{eq:RQpub}
\end{align}

Informally: the CODATA/NIST measured value differs from the pure source calculation by a small, precisely localized inverse-alpha contribution. The golden-ratio equation, the heat-kernel closure, and the numerical fixed-point solve all point to the same address for the difference: the same-scheme low-energy hadronic transport.

Why this step is needed: this accounting prevents the hadronic contribution from being hidden inside the final number. It shows the source anchor, the calculated transport, the required transport, and the residual in the same units.

\section{Step 18: The OPH Output and the Required Hadronic Correction}

The historical printed run of the fixed-point calculation emits
\begin{equation}
A_{\mathrm{calc}}(P_{\mathrm{root}})
=
\alphainv_{\mathrm{root}}
=
\decnum{136.994835164621649457949994585787193262029}.
\label{eq:Acalcfp}
\end{equation}
A converged precision-100 rerun supersedes this printed tail beyond digit~9;
the certified value is \(\alphainv_{\mathrm{root}}=136.994835177413\ldots\)
(ledger row CL-6, closed; the printed-pair identity is pinned by a
CI test in the code release \cite{oph-code}). The remaining
arithmetic in this section keeps the printed tail so that the quoted display
packet stays internally consistent as a historical record.
At the CODATA-derived comparison pixel, the finite-screen unified gauge-width contribution is
\[
\alpha_U(\Ppub)=0.041124336195630495.
\]
Combining that comparison-pixel gauge width with the source/root value gives the mixed-provenance
no-hadron diagnostic
\begin{equation}
A_{\alpha_U}^{\mathrm{fp}}
=\alpha_{\mathrm{root}}^{-1}+\alpha_U(\Ppub)
=137.035959500817279952949994585787\ldots,
\label{eq:barealphaUaddition}
\end{equation}
or
\begin{equation}
\alpha_{\mathrm{src}+U}^{\mathrm{fp}}
=\bigl(A_{\alpha_U}^{\mathrm{fp}}\bigr)^{-1}
=0.00729735467714250592971405837329\ldots.
\label{eq:barealphaUalpha}
\end{equation}
This \(A_{\alpha_U}^{\mathrm{fp}}\) value is comparison bookkeeping, not a source-only prediction:
\(\alpha_{\mathrm{root}}^{-1}\) is evaluated on the source root, while \(\alpha_U(\Ppub)\) inherits
the measured endpoint through \(\Ppub\). It is a mixed-provenance display packet (inner value
evaluated at \(P_{\mathrm{fwd}}\) plus \(\alpha_U\) evaluated at the SL-3 pixel); it is not a fixed
point of any single declared map, and it stays on record as a display packet only. The certified
self-consistent gauge-width fixed point is \(\alpha^{-1}=137.035660136946577\ldots\) (ledger row
CL-2).

Compared with the CODATA/NIST central measured endpoint,
\begin{equation}
A_{\mathrm C}=137.035999177,
\qquad
\sigma_A=0.000000021,
\end{equation}
the CODATA-minus-diagnostic bookkeeping gap is
\begin{equation}
\Dhadreq^{\mathrm{fp}}
=
A_{\mathrm C}-A_{\alpha_U}^{\mathrm{fp}}
=
\decnum{0.000039676182720047050005414213}\ldots.
\label{eq:Dempfp}
\end{equation}
Equivalently,
\begin{equation}
\Dhadreq^{\mathrm{fp}}
=0.0000396762\pm0.000000021,
\qquad
\frac{\Dhadreq^{\mathrm{fp}}}{A_{\mathrm C}}
=2.89531\times10^{-7}.
\label{eq:Dhadrequncertainty}
\end{equation}
This is the inverse-alpha gap between the mixed diagnostic and the CODATA/NIST central value. Its
physical address is the low-energy QCD/hadronic vacuum-polarization and same-scheme transport of the
Ward-projected electromagnetic current, but the subtraction does not calculate that transport.
The \(2.9\times10^{-7}\) ratio describes the display packet only. The ledger residual CL-2 is
measured against the certified self-consistent gauge-width fixed point
\(\alpha^{-1}=137.035660136946577\ldots\) and is \(2.5\times10^{-6}\) relative, about
\(1.6\times10^{4}\) measurement sigma; that residual supersedes the \(2.9\times10^{-7}\) figure as
the closure status of the loop.

For bookkeeping, the full root-to-CODATA difference is
\begin{align}
\Dcalib^{\mathrm{fp}}
&=
A_{\mathrm C}-A_{\mathrm{calc}}(P_{\mathrm{root}})
\nonumber\\
&=
\decnum{0.041164012378350542050005414212806737971}
\nonumber\\
&=
\alpha_U(\Ppub)+\Dhadreq^{\mathrm{fp}}.
\label{eq:rootcodatafp}
\end{align}
This larger difference is bookkeeping: it contains the comparison-pixel unified gauge-width
contribution and the CODATA-minus-diagnostic gap.

Equivalently, the CODATA/NIST central value can be written as
\begin{equation}
A_{\mathrm C}
=
\alpha_{\mathrm{root}}^{-1}+\alpha_U(\Ppub)\,C_{24,Q},
\qquad
C_{24,Q}=1.0009647859732326253849511140702475\ldots.
\label{eq:c24qfactor}
\end{equation}
The factor \(C_{24,Q}\) is a back-solved calibration: compact comparison
accounting, solved from the measured endpoint, for the finite same-scheme hadronic endpoint
transport needed for source-only closure. Equivalently,
\(\Dhadreq^{\mathrm{fp}}=\alpha_U(\Ppub)(C_{24,Q}-1)\). The factor is not a source-emitted
hadronic/QCD payload.

Here \(A_{\mathrm C}\) is the CODATA/NIST 2022 inverse fine-structure constant \cite{nist-value,nist-2022}. The source value \(A_{\mathrm{calc}}(P_{\mathrm{root}})=\alpha_{\mathrm{root}}^{-1}\) is a reproducible source-root output \cite{oph-code}; \(A_{\alpha_U}^{\mathrm{fp}}\) mixes that root with \(\alpha_U(\Ppub)\). Equation~\eqref{eq:Dempfp} is the measured-comparison residual in inverse-alpha units.

Informally: the mixed comparison diagnostic is \(137.035959500817\ldots\), a display packet; the
certified self-consistent gauge-width fixed point is \(137.035660136946577\ldots\) (row CL-2,
\(2.5\times10^{-6}\) relative); and the CODATA/NIST measured value is \(137.035999177(21)\). The
remaining difference is the size required of the missing same-scheme endpoint contribution; the
open source task is to compute that contribution from an OPH-QCD spectral backend.

This marks the calculation boundary: additional digits require a source-derived hadronic spectral calculation.

The endpoint decomposition is
\begin{equation}
\ATh(P)
=
\AZ(P)+\Dlep(P)+\Dhad(P)+\DEW(P).
\label{eq:ATsplit}
\end{equation}
The calculated transport contains the exact one-loop charged-lepton contribution and a structured quark-screening continuation. The remaining pure source object can be written compactly as
\begin{equation}
\Rq(P)
=
\Dhad(P)+\DEW(P)
-\bigl[\Dcalc(P)-\Dlep(P)\bigr],
\label{eq:RQmeaning}
\end{equation}
where the right side is understood in the same renormalization and endpoint scheme as \(\AZ(P)\).

In the screening notation used by the endpoint calculation, let
\begin{equation}
x(P)=\frac{\Nc\alpha_3(m_Z;P)}{\pi}.
\label{eq:xqcd}
\end{equation}
At \(P=\Ppub\),
\begin{align}
S_{\mathrm{required}}
&=
\decnum{0.89540013264765879780580028318167064130770986481229164844830591011545242521273086888107864576}
\label{eq:Sreq}\\
c_Q
&=
\decnum{0.65802575992715543563823017023236005042492009907058639608566007832470711257322011342309305088}
\label{eq:cQ}
\end{align}
where
\begin{equation}
S_{\mathrm{required}}=1-x+c_Qx^2.
\label{eq:cQdef}
\end{equation}

Informally: the needed hadronic and endpoint correction can be summarized as a small second-order screening coefficient. A pure source proof has to emit this coefficient from a Ward-projected hadronic spectral measure.

Why this step is needed: hadrons are not pointlike free quarks at low energy. The coefficient is a compact way to record what the empirical hadronic spectral function contributes in this endpoint convention.

\section{Step 19: Why Raw PDG \(\Delta\alpha_{\mathrm{had}}^{(5)}(M_Z)\) Is a Different Quantity}

The direct PDG/CERN diagnostic uses
\begin{align}
\AZ(\Ppub)&=
\decnum{128.30796547328624820996110874175671618724547618036535646005342169635117784168285644}\\
\Dlep(\Ppub)&=
\decnum{4.3093978664522040271317438975344894018487156605576773194711528089665680313257906466129}\\
\Delta\alpha_{\mathrm{had}}^{(5)}(M_Z)&=0.02761.
\end{align}
It forms
\begin{equation}
A_L=\AZ+\Dlep,
\qquad
A_{\mathrm{PDGdiag}}=\frac{A_L}{1-\Delta\alpha_{\mathrm{had}}^{(5)}(M_Z)}.
\label{eq:pdgdiag}
\end{equation}
Numerically,
\begin{equation}
A_{\mathrm{PDGdiag}}
=
\decnum{136.382895072695577121415124218977165118002233508081154453999500720202537945689123794581289}
\label{eq:pdgdiagval}
\end{equation}
This differs from \(137.035999177\). Holding this raw hadronic denominator shift fixed, the target would require
\begin{equation}
\Delta\alpha_{\mathrm{had,req}}
=
\decnum{0.032244343557887288822268499236957947422007347612226081506593089221871380926117950141438549}
\label{eq:pdgreq}
\end{equation}
or a same-scheme source-anchor bridge of
\begin{equation}
\decnum{0.6350718999845777629071473607087944109058081590769662204754254946822541269913529133871}
\label{eq:bridgegap}
\end{equation}
inverse-alpha units.

Informally: the raw electroweak-review hadronic running number is a useful diagnostic. The OPH endpoint needs the same electromagnetic current and the same inverse-alpha endpoint convention used by \(\AZ(P)\). Mixing the two conventions moves the answer by a visible amount.

Why this step is needed: a common mistake is to insert a published \(\Delta\alpha_{\mathrm{had}}^{(5)}(M_Z)\) number as though it were the OPH endpoint contribution. This section shows the numerical consequence of that convention mismatch.

\section{Step 20: The Measured Hadronic Spectral Input}

The standard measured hadronic vacuum-polarization relation uses
\begin{equation}
\Delta\alpha_{\mathrm{had}}(q^2)
=
-\frac{\alpha q^2}{3\pi}\,
\mathrm{P.V.}\!\int \frac{R(s)}{s(s-q^2)}\dd s,
\label{eq:dispersion}
\end{equation}
where \(R(s)\) is the measured ratio of the bare \(e^+e^-\to\mathrm{hadrons}\) cross section to the pointlike \(e^+e^-\to\mu^+\mu^-\) cross section.

The independently documented measured quantity in the standard electroweak literature is usually quoted as \(\Delta\alpha_{\mathrm{had}}^{(5)}(M_Z^2)\). It is a dimensionless change in the running electromagnetic coupling at the \(Z\)-boson mass, not the inverse-alpha correction \(\Dhadreq^{\mathrm{fp}}\) used in Eq.~\eqref{eq:Dempfp}. For example, Davier, Hoecker, Malaescu, and Zhang report
\[
\Delta\alpha_{\mathrm{had}}^{(5)}(M_Z^2)=(275.7\pm1.0)\times10^{-4}
\]
from \(e^+e^-\)-based data \cite{dhmz}. A perturbative update by Erler and Ferro-Hernandez gives
\[
\Delta\alpha_{\mathrm{had}}^{(5)}(M_Z^2)
=\left(276.29\pm0.38\pm0.62\right)\times10^{-4}
\]
when low-energy cross-section data are used as input \cite{erler-ferro}.

The comparison residual in Eq.~\eqref{eq:Dempfp} is different. It is expressed in inverse-alpha units and in the endpoint convention used by \(\AZ(P)\). The standard published hadronic-running values document the physics source of the correction. They do not by themselves give the OPH same-scheme inverse-alpha correction to arbitrary precision. The community does publish the data machinery needed for this kind of calculation. Jegerlehner's alphaQED package provides hadronic running routines, covariance data, and an integration routine for custom kernels \cite{alphaqed}. That is the right empirical route once the OPH endpoint kernel and finite same-scheme remainder are fixed.

In the OPH endpoint convention this is represented as a same-current spectral functional,
\begin{equation}
\Dhad(P)
=
\frac{m_Z(P)^2}{3\pi}
\int \frac{\rho_Q(s;P)}{s[s+m_Z(P)^2]}\dd s,
\label{eq:ophhadron}
\end{equation}
plus the same-scheme finite remainder needed by Eq.~\eqref{eq:ATsplit}.

Informally: the measured hadronic input supplies the electromagnetic spectral information required for the endpoint. The single published number \(0.0276\) is one weighted integral of that spectral information. The OPH endpoint needs a different weighted integral plus the same-scheme finite remainder. A pure source theorem requires the same spectral object from OPH source data.

Why this step is needed: the dispersion relation explains why measured \(e^+e^-\to\mathrm{hadrons}\) data are the right empirical input. They measure the electromagnetic spectral function that the endpoint transport requires.

\section{No-Go: The 24-Register Scale Does Not Determine the Hadronic Residual}

\begin{theorem}[Coarse 24-register data do not determine the Thomson residual]
The specified OPH source objects \(P\), \(\alpha_U(P)\), \(N_c=3\), \(N_g=3\), and the
24-slot repair-register count do not determine the exact low-energy Thomson residual.
\end{theorem}

\begin{proof}
The hadronic endpoint contribution depends on the Ward-projected spectral functional
\begin{equation}
\Delta_{\mathrm{had}}(P)
=
\frac{m_Z(P)^2}{3\pi}
\int_0^\infty
\frac{\rho_Q(s;P)}{s[s+m_Z(P)^2]}\dd s
\label{eq:hadronkernelnogo}
\end{equation}
up to the same-scheme finite endpoint remainder. Let
\[
K_P(s)=\frac{m_Z(P)^2}{3\pi s[s+m_Z(P)^2]}.
\]
This kernel is not constant. Two positive spectral measures can therefore have the same coarse
register count, the same leading normalization, and the same electroweak branch data while giving
different residuals. For example, with \(s_1\ne s_2\),
\[
\rho_1(s)=M\delta(s-s_1),
\qquad
\rho_2(s)=M\delta(s-s_2)
\]
have equal total mass, but
\[
\int K_P(s)\rho_1(s)\dd s=M K_P(s_1),
\qquad
\int K_P(s)\rho_2(s)\dd s=M K_P(s_2),
\]
which are generically unequal. Thus the 24-register scale explains the size of the correction but
does not select the spectral measure that fixes its exact value.
\end{proof}

The hadronic precision audit sharpens the no-go statement. The
\(\rho_Q(s;P)\) object used here is the two-current marginal of a larger OPH-QCD
hadronic backend: a source QCD quotient ensemble, source parameter map, Ward-normalized current
ledger, Stieltjes/Jacobi export, same-scheme remainder, and systematics/no-target-leak receipt
bundle. That marginal can close the running-\(\alpha\)/HVP transport only after it is emitted by
the backend. It cannot be recycled as a full hadronic precision source, because HLbL and rare
decay long-distance amplitudes require four-current and transition spectral data from the same
source law.

\begin{theorem}[Tautological 24-Jacobi realization]
For any positive residual \(R>0\), there is an exact 24-point Jacobi/Stieltjes representation that
reproduces \(R\). Such a representation is non-predictive unless its nodes, weights,
normalization, and finite remainder are emitted by OPH source rules before comparison with the
Thomson endpoint.
\end{theorem}

\begin{proof}
Choose any positive spectral nodes \(s_1,\ldots,s_{24}>0\) and any positive normalized weights
\(u_j>0\) with \(\sum_j u_j=1\). The finite inverse Stieltjes/Favard construction gives a
positive \(24\times24\) Jacobi matrix \(J_{24,Q}\) whose \(e_1\)-spectral measure is
\[
\dd\mu_0(s)=\sum_{j=1}^{24}u_j\delta(s-s_j).
\]
Equivalently,
\[
e_1^\top f(J_{24,Q})e_1=\sum_{j=1}^{24}u_j f(s_j)
\]
for every function \(f\) on the spectrum. Define
\[
D(J_{24,Q};m)
=
e_1^\top\!\left[J_{24,Q}^{-1}-(J_{24,Q}+m^2I)^{-1}\right]e_1.
\]
Since \(s_j>0\),
\[
D(J_{24,Q};m)=
\sum_{j=1}^{24}u_j\left(\frac1{s_j}-\frac1{s_j+m^2}\right)>0.
\]
For any same-scheme finite remainder \(\Xi_Q<R\), set
\[
\omega_Q=\frac{3\pi(R-\Xi_Q)}{D(J_{24,Q};m)}.
\]
Then
\[
\frac{\omega_Q}{3\pi}
e_1^\top\!\left[J_{24,Q}^{-1}-(J_{24,Q}+m^2I)^{-1}\right]e_1
+\Xi_Q
=R.
\]
The construction reproduces the chosen residual exactly, but the freedom in \(s_j\), \(u_j\),
\(\omega_Q\), and \(\Xi_Q\) means it is a back-solving theorem rather than a source derivation.
\end{proof}

The source-side object needed for a final OPH endpoint theorem is therefore
\begin{equation}
\Rq(P)
=
\frac{\omega_Q(P)}{3\pi}
e_1^\top\!\left[
J_{24,Q}(P)^{-1}
-
\bigl(J_{24,Q}(P)+m_Z(P)^2I\bigr)^{-1}
\right]e_1
+\Xi_Q(P),
\label{eq:sourceJ24Q}
\end{equation}
or equivalently the source-emitted spectral density \(\rho_Q(s;P)\) plus the same-scheme finite
remainder. The required dependency graph must have no path from the measured Thomson endpoint into
\(J_{24,Q}\), \(\omega_Q\), \(\Xi_Q\), or \(\rho_Q\). The historical v2 endpoint target
and its OpenTimestamps files remain an auditable timestamp record, but they do not define an
active blind test.  The V1 execution embedded target constants, was directed by a
session with target access, used a different \(P\) from the target, confused the closure residual
with the required total payload, and supplied a sampled min/max envelope rather than a certified
interval.  A corrected contract must freeze one coordinate schema, require a source-emitted
function over the same \(P\), certify quadrature and interval bounds, and score the emitted
artifact separately.  No promotion follows from v2 or V1.

\section{Step 21: The Conditional Pure Source Theorem}

\begin{theorem}[OPH fine-structure endpoint, conditional pure source form]
Assume:
\begin{enumerate}[label=(\roman*)]
\item the OPH overlap-consistency branch emits the source map \(P\mapsto \AZ(P)\) by Eqs.~\eqref{eq:MU}--\eqref{eq:AZ};
\item the Ward-projected \(\U(1)_Q\) transport theorem emits a same-scheme endpoint map
\[
\ATh(P)=\AZ(P)+\Dlep(P)+\Dhad(P)+\DEW(P);
\]
\item \(G(P)=\varphi+\sqrt\pi/\ATh(P)\) is a self-map and a contraction on the physical pixel interval \(I\);
\item the interval image contains the root and the residual bound is certified.
\end{enumerate}
Then there is a unique \(P_\star\in I\) satisfying \(G(P_\star)=P_\star\), and the fine-structure constant on that branch is
\[
\alpha(0)=\frac{1}{\ATh(P_\star)}=\frac{P_\star-\varphi}{\sqrt\pi}.
\]
\end{theorem}

\begin{proof}
By assumption (iii), \(G:I\to I\) is a contraction. Banach's fixed-point theorem gives a unique \(P_\star\in I\) such that \(G(P_\star)=P_\star\). Equation~\eqref{eq:G} gives
\[
P_\star-\varphi=\frac{\sqrt\pi}{\ATh(P_\star)}.
\]
Dividing by \(\sqrt\pi\) gives
\[
\frac{P_\star-\varphi}{\sqrt\pi}=\frac{1}{\ATh(P_\star)}.
\]
The left side is \(\alpha_{\mathrm{ext}}(P_\star)\), and the right side is \(\alpha_{\mathrm{in}}(P_\star)\). Their common value is the Thomson-limit electromagnetic coupling \(\alpha(0)\).
\end{proof}

Informally: a source-derived hadronic spectral endpoint map would close the last open transport step. With the interval proof included, the fine-structure constant is the unique fixed point of the full source map.

Uniqueness lemma L1 of the project consistency stack
states the interval form of assumptions
(iii)--(iv): on any interval \(I\) where interval-arithmetic evaluation
certifies \(G(I)\subseteq I\) together with a derivative bound
\(|G'|\le L<1\), the Banach fixed-point theorem gives both existence and
uniqueness of the fixed point in \(I\). L1 is discharged on an explicit
interval. The interval contraction certificate shipped with
the code release \cite{oph-code},
proves \(G(I)\subseteq\operatorname{int}(I)\) and \(L\le0.0724\) for both
readout modes by a direct Banach argument in mean-value (centered) form, with
\texttt{mpmath.iv} outward rounding on every elementary operation and with the
\(\SU(2)\)/\(\SU(3)\) edge-sum tails bounded by geometric majorants, so the
enclosure covers the infinite-cutoff sums. The certified unique fixed points
are \(\alpha^{-1}=136.994835177413\ldots\) for the source map (enclosure width
\(7.2\times10^{-24}\)) and \(\alpha^{-1}=137.035660136946577\ldots\) for the
gauge-width map. The certificate certifies the declared numerical map; it does
not relate either fixed point to the measured constant, and the stage-3
landing verdict of the basin-then-contract protocol is unchanged (outside the
SL-3 basin, rows CL-1 and CL-2). The global at-most-one statement is
discharged on the declared physical domain (\(\alpha^{-1}\in[100,200]\)):
\(\sup|g'|<1\) on all 256 certified pieces, both readout maps, zero
exceptional set
(domain-global uniqueness certificate in the same release), so
each map has exactly one fixed point on the domain.

The release gate for promoting this row from empirical endpoint comparison to source-only theorem is
explicit. The packet must contain
\[
\texttt{WARD\_Q\_SPECTRAL\_MEASURE\_RECEIPT},\quad
\texttt{J24Q\_SOURCE\_JACOBI\_RECEIPT},\quad
\texttt{OMEGAQ\_NORMALIZATION\_RECEIPT},
\]
\[
\texttt{XI\_SAME\_SCHEME\_REMAINDER\_RECEIPT},\quad
\texttt{NO\_THOMSON\_TARGET\_LEAK\_DAG},\quad
\texttt{FULL\_PIXEL\_CONTRACTION\_INTERVAL},
\]
and \(\texttt{ALPHA\_PREDICTION\_BEATS\_DIRECT\_MEASUREMENT}\) under the direct-measurement
interval adopted for that release. For the present comparison standard, the direct rubidium recoil
result of Morel, Yao, Cladé, and Guellati-Khélifa reports
\(\alpha^{-1}=137.035999206(11)\) with 81 parts-per-trillion relative precision
\cite{morel2020}. The operational receipt is therefore
\[
\operatorname{rad}(\alpha_{\mathrm{pred}}^{-1})<1.1\times10^{-8}
\]
in inverse-alpha units, unless a later release updates the adopted direct-measurement interval.

\section{Step 22: The CODATA/NIST Comparison Endpoint}

With the CODATA/NIST measured endpoint value
\begin{equation}
\ATh(\Ppub)=A_{\mathrm C}=137.035999177(21),
\end{equation}
Eq.~\eqref{eq:G} gives the corresponding comparison pixel:
\begin{align}
\boxed{
\Ppub
=
\decnum{1.6309682094039593248792798477826489413359828516279250606661507533907793398933432}
}
\\
\boxed{
\alpha(0)
=
\decnum{0.0072973525643314250302457952646916832280660213133653604957798803819933561573639928}
}
\\
\boxed{
\alphainv(0)=137.035999177(21).
}
\end{align}

Informally: the endpoint number is the same number NIST/CODATA reports. The derivation exposes every OPH-side step and identifies where the low-energy hadronic calculation enters. The missing hadronic spectral calculation is work in progress.

\section{Reproducibility}

The executable derivation lives in the public repository under
\url{https://github.com/FloatingPragma/observer-patch-holography/tree/main/code/P_derivation}.
The main human-facing CLI is
\begin{verbatim}
cd reverse-engineering-reality/code/P_derivation
python3 derive_p.py --color always
\end{verbatim}
To print only the pure source value:
\begin{verbatim}
python3 derive_p.py --no-hadron-closure
\end{verbatim}
To emit machine-readable output:
\begin{verbatim}
python3 derive_p.py --json --output runtime/report.json
\end{verbatim}
To run the raw PDG/CERN diagnostic discussed in Section~19:
\begin{verbatim}
python3 fine_structure_fixed_point_demo.py --compare-alpha-inv 137.035999177
\end{verbatim}

Informally: the command line separates the pure source value, the CODATA/NIST comparison value, and the diagnostic values. The mathematical paper uses the same separation.

\section{Checklist of Non-Omitted Steps}

\begin{longtable}{@{}L{0.18\textwidth}L{0.44\textwidth}L{0.28\textwidth}@{}}
\toprule
Step & Mathematical object & Status \\
\midrule
\endhead
Overlap rule & Eq.~\eqref{eq:overlap} & OPH starting structure. \\
Pixel variable & \(P=a_{\mathrm{cell}}/\ell_\star^2\) & Defined. \\
Golden balance & \(\varphi=(1+\sqrt5)/2\) & Defined. \\
Boundary width & \(\alpha_{\mathrm{ext}}=(P-\varphi)/\sqrt\pi\) & Defined. \\
Inside readout & \(\alpha_{\mathrm{in}}=1/\ATh(P)\) & Defined. \\
Closure & \(P=\varphi+\sqrt\pi/\ATh(P)\) & Defined. \\
Source scale & \(M_U=E_P\e^{-2\pi}P^{1/6}\) & Defined and calculated. \\
Cell scale & \(E_{\mathrm{cell}}=E_P/\sqrt P\) & Defined and calculated. \\
Transmutation & \(v=E_{\mathrm{cell}}\exp[-2\pi/(\beta_{\mathrm{EW}}\alpha_U)]\) & Defined and calculated. \\
Gauge running & Eq.~\eqref{eq:running} & Defined and calculated. \\
\(Z\)-scale self-consistency & Eq.~\eqref{eq:mZfixed} & Defined and calculated. \\
Heat-kernel closure & Eq.~\eqref{eq:pixelclosure} & Defined and calculated. \\
Electroweak anchor & \(\AZ=\alpha_{\mathrm{em}}^{-1}(m_Z^2;P)\) & Defined and calculated. \\
Charged spectrum & Eqs.~\eqref{eq:quarkmasses} and \eqref{eq:lepmasses} & Frozen phenomenological continuation; not an OPH-derived mass prediction. \\
Fermion kernel & Eqs.~\eqref{eq:kernelint} and \eqref{eq:kernelclosed} & Exact one-loop kernel. \\
Calculated endpoint & \(A_{\mathrm{calc}}=\AZ+\Dcalc\) & Raw no-hadron diagnostic conditional on the supplied charged-spectrum continuation. \\
Bare-\(\alpha_U\) closure status & \(A_{\alpha_U}^{\mathrm{fp}}\), Eq.~\eqref{eq:barealphaUaddition} & Mixed-provenance display packet; the CL-2 closure status is the certified gauge-width fixed point \(137.035660136946577\ldots\); not a derivation of \(\alpha\). \\
Required hadronic correction & \(\Dhadreq^{\mathrm{fp}}\), Eq.~\eqref{eq:Dempfp} & CODATA/NIST comparison residual attributed to the missing same-scheme hadronic endpoint transport. \\
Hadronic/scheme residual & \(\Rq=\ATh-A_{\mathrm{calc}}\) & Isolated endpoint contribution. \\
CODATA/NIST endpoint & \(\ATh=137.035999177(21)\) & Measured endpoint comparison value. \\
Two-\(P\) provenance & \(P_{\mathrm{root}}\) versus \(\Ppub\) & Source-branch root and measured-endpoint comparison pixel are distinct; \(\Ppub\) inherits CODATA/NIST by definition. \\
Raw PDG diagnostic & Eq.~\eqref{eq:pdgdiag} & Comparison diagnostic in a different endpoint convention. \\
Interval theorem & Banach/contraction certificate for full \(G\) & Discharged for the declared numerical map (shipped certificate, both readout modes); still required for the full source map with the hadronic term. \\
\bottomrule
\end{longtable}

\section{Conclusion}

The fine-structure constant is the coupling that makes one OPH screen cell internally consistent. The outside reading says that the cell is displaced from the golden-ratio balance by \((P-\varphi)/\sqrt\pi\). The inside reading sends the same \(P\) through the source map, gauge closure, electroweak anchor, and Ward-projected electromagnetic transport, producing \(1/\ATh(P)\). The physical cell is the fixed point where those two readings agree. The construction carries zero continuous dials, each declared map has a machine-certified unique fixed point on the declared domain, and the open hadronic step sits inside a standing falsification program with frozen, externally timestamped targets, so every residual quoted below is a measurement of the principle set rather than a fit.

The pure source calculation computes the root witness
\[
\alphainv_{\mathrm{root}}
=
\decnum{136.994835177412937295289429464436902857561206151035393184502},
\]
the interval-certified unique fixed point of the declared source map
(enclosure width \(7.2\times10^{-24}\), shipped interval certificate).
The CODATA/NIST measured endpoint value is
\[
\alphainv(0)=137.035999177(21),
\]
with
\[
P=1.630968209403959324879279847782648941\ldots.
\]
Combining the source root with the finite-screen unified gauge-width contribution evaluated at the
CODATA-derived comparison pixel gives the mixed-provenance no-hadron diagnostic
\[
A_{\alpha_U}^{\mathrm{fp}}=137.035959500817279952949994585787\ldots.
\]
That printed value combines the superseded unconverged root tail with \(\alpha_U(\Ppub)\); it is a
display packet (inner value at \(P_{\mathrm{fwd}}\) plus \(\alpha_U\) at the SL-3 pixel), not a
fixed point of any single declared map, and it stays on record as a display packet only. The
certified self-consistent gauge-width fixed point is
\(\alpha^{-1}=137.035660136946577\ldots\), \(2.5\times10^{-6}\) below the measured value, about
\(1.6\times10^{4}\) measurement sigma (row CL-2).
The remaining gap is CODATA-minus-diagnostic bookkeeping, not an emitted low-energy hadronic
calculation. The pure source theorem requires the Ward-projected hadronic spectral measure, the
same-scheme bridge, and the interval-certificate step.

\section*{Declarations}

\textbf{Funding declaration.} No funding was received for this work.

\textbf{Consent to Participate declaration.} Not applicable.

\textbf{Consent to Publish declaration.} Not applicable.

\textbf{Author Contribution declaration.} Bernhard M\"uller conceived the OPH fine-structure fixed-point formulation, prepared the mathematical derivation, wrote and checked the numerical calculation, wrote the manuscript, and prepared the reproducibility materials. Alexander Osika contributed OPH research-program framing, manuscript review, and approval of the final manuscript.

\textbf{Data Availability declaration.} No new experimental data were generated. The manuscript cites the 2022 CODATA/NIST inverse fine-structure constant \cite{nist-value,nist-2022}, public hadronic cross-section and vacuum-polarization references \cite{pdg-xsect,dhmz,knt19}, and the OPH source materials cited below. The TeX source, OPH paper sources, and executable derivation code are available in the public GitHub repository:
\url{https://github.com/FloatingPragma/observer-patch-holography}.

\textbf{Code Availability declaration.} The fixed-point code and runtime artifacts are available at
\url{https://github.com/FloatingPragma/observer-patch-holography/tree/main/code/P_derivation}.

\textbf{Ethics declaration.} Not applicable.

\textbf{Competing Interests declaration.} The authors declare no competing interests.

\begin{thebibliography}{99}

\bibitem{oph-synthesis}
B.\ M\"uller, A.\ Osika, M.\ Poneder, K.\ Xue, B.\ Cassie, P.\ Nguyen, M.\ A.\ Visser, K.\ A.\ Anirudha, D.\ Matscheko, and J.\ Hill,
\emph{Observers Are All You Need}.
GitHub source:
\url{https://github.com/FloatingPragma/observer-patch-holography/blob/main/paper/observers_are_all_you_need.tex}.
GitHub PDF:
\url{https://github.com/FloatingPragma/observer-patch-holography/blob/main/paper/observers_are_all_you_need.pdf}.

\bibitem{oph-smgr}
B.\ M\"uller, A.\ Osika, M.\ Poneder, K.\ Xue, P.\ Nguyen, M.\ A.\ Visser, and D.\ Matscheko,
\emph{Recovering Relativity and the Standard Model from Observer Overlap Consistency}.
GitHub source:
\url{https://github.com/FloatingPragma/observer-patch-holography/blob/main/paper/recovering_relativity_and_standard_model_structure_from_observer_overlap_consistency_compact.tex}.
GitHub PDF:
\url{https://github.com/FloatingPragma/observer-patch-holography/blob/main/paper/recovering_relativity_and_standard_model_structure_from_observer_overlap_consistency_compact.pdf}.

\bibitem{oph-particle}
B.\ M\"uller, A.\ Osika, M.\ Poneder, K.\ Xue, M.\ A.\ Visser, and D.\ Matscheko,
\emph{Deriving the Particle Zoo from Observer Consistency}.
GitHub source:
\url{https://github.com/FloatingPragma/observer-patch-holography/blob/main/paper/deriving_the_particle_zoo_from_observer_consistency.tex}.
GitHub PDF:
\url{https://github.com/FloatingPragma/observer-patch-holography/blob/main/paper/deriving_the_particle_zoo_from_observer_consistency.pdf}.

\bibitem{oph-code}
OPH executable derivation code,
\url{https://github.com/FloatingPragma/observer-patch-holography/tree/main/code/P_derivation}.
Primary files:
\path{paper_math.py},
\path{derive_p.py},
\path{fine_structure_fixed_point_demo.py}.

\bibitem{oph-hadron-policy}
OPH hadron data policy,
\url{https://github.com/FloatingPragma/observer-patch-holography/blob/main/docs/HADRON.md}.

\bibitem{jentschura-nandori}
U.\ D.\ Jentschura and I.\ Nandori,
\emph{Attempts at a determination of the fine-structure constant from first principles: A brief historical overview},
Eur.\ Phys.\ J.\ H \textbf{39}, 591--613 (2014).
\url{https://arxiv.org/abs/1411.4673}.

\bibitem{heyrovska-narayan}
R.\ Heyrovska and S.\ Narayan,
\emph{Fine-structure Constant, Anomalous Magnetic Moment, Relativity Factor and the Golden Ratio that Divides the Bohr Radius}.
\url{https://arxiv.org/abs/physics/0509207}.

\bibitem{singh}
T.\ P.\ Singh,
\emph{Quantum theory without classical time: octonions, and a theoretical derivation of the fine structure constant \(1/137\)},
Int.\ J.\ Mod.\ Phys.\ D \textbf{30}, 2142010 (2021).
\url{https://arxiv.org/abs/2110.07548}.

\bibitem{ciborowski}
J.\ Ciborowski,
\emph{Bi-Constructible pattern of weak and flavour mixing: implications for electroweak coupling constants}.
\url{https://arxiv.org/abs/2508.00030}.

\bibitem{chekanov-kjellerstrand}
S.\ V.\ Chekanov and H.\ Kjellerstrand,
\emph{Discovering the Underlying Analytic Structure Within Standard Model Constants Using Artificial Intelligence},
Particles \textbf{8}, 95 (2025).
\url{https://arxiv.org/abs/2507.00225}.

\bibitem{liu-low-yin}
Q.\ Liu, I.\ Low, and Z.\ Yin,
\emph{A Quantum Computational Determination of the Weak Mixing Angle in the Standard Model}.
\url{https://arxiv.org/abs/2509.18251}.

\bibitem{dhmz}
M.\ Davier, A.\ Hoecker, B.\ Malaescu, and Z.\ Zhang,
\emph{Reevaluation of the Hadronic Contributions to the Muon \(g-2\) and to \(\alpha(M_Z)\)},
Eur.\ Phys.\ J.\ C \textbf{71}, 1515 (2011).
\url{https://arxiv.org/abs/1010.4180}.

\bibitem{knt19}
A.\ Keshavarzi, D.\ Nomura, and T.\ Teubner,
\emph{The \(g-2\) of charged leptons, \(\alpha(M_Z^2)\) and the hyperfine splitting of muonium},
Phys.\ Rev.\ D \textbf{101}, 014029 (2020).
\url{https://arxiv.org/abs/1911.00367}.

\bibitem{morel2020}
L.\ Morel, Z.\ Yao, P.\ Cladé, and S.\ Guellati-Khélifa,
\emph{Determination of the fine-structure constant with an accuracy of 81 parts per trillion},
Nature \textbf{588}, 61--65 (2020).
\url{https://doi.org/10.1038/s41586-020-2964-7}.

\bibitem{erler-ferro}
J.\ Erler and R.\ Ferro-Hernandez,
\emph{Perturbative contributions to \(\Delta\alpha^{(5)}(M_Z^2)\)}.
\url{https://arxiv.org/abs/2308.05740}.

\bibitem{pdg-xsect}
Particle Data Group,
\emph{Data files and plots of cross-sections and related quantities in the Review of Particle Physics}.
\url{https://pdg.lbl.gov/current/xsect/}.

\bibitem{alphaqed}
F.\ Jegerlehner,
\emph{Software packages alphaQED and pQCDAdler}.
\url{https://people.physik.hu-berlin.de/~fjeger/software.html}.

\bibitem{blog}
B.\ M\"uller,
\emph{Fine-Structure Constant from First Principles},
Floating Pragma Blog.
\url{https://blog.floatingpragma.io/fine-structure-constant-from-first-principles}.

\bibitem{nist-value}
NIST,
\emph{CODATA Value: inverse fine-structure constant},
2022 CODATA recommended value.
\url{https://physics.nist.gov/cgi-bin/cuu/Value?eqalphinv=}.

\bibitem{nist-2022}
P.\ J.\ Mohr, E.\ Tiesinga, D.\ B.\ Newell, and B.\ N.\ Taylor,
\emph{CODATA Recommended Values of the Fundamental Physical Constants: 2022}.
NIST publication page:
\url{https://www.nist.gov/publications/codata-recommended-values-fundamental-physical-constants-2022}.

\bibitem{nist-wall}
NIST/CODATA,
\emph{CODATA Recommended Values of the Fundamental Physical Constants: 2022},
wall chart.
\url{https://physics.nist.gov/cuu/pdf/wall_2022.pdf}.

\end{thebibliography}

\end{document}
